%-------------------------
% Resume in Latex
% Author : Harshibar
% Based off of: https://github.com/jakeryang/resume
% License : MIT
%------------------------

\documentclass[letterpaper,11pt]{article}

\usepackage{latexsym}
\usepackage[empty]{fullpage}
\usepackage{titlesec}
\usepackage{marvosym}
\usepackage[usenames,dvipsnames]{color}
\usepackage{verbatim}
\usepackage{enumitem}
\usepackage[hidelinks]{hyperref}
\usepackage{fancyhdr}
\usepackage{fontspec}
\defaultfontfeatures{Ligatures=TeX}
\setmainfont{Arial}
\setsansfont{Arial}
\setmonofont{Noto Sans Mono}[Scale=0.90]
\usepackage[main=russian,english]{babel}
\usepackage{tabularx}
% only for pdflatex
% \input{glyphtounicode}

% fixed width font is configured through fontspec

% light-grey
\definecolor{light-grey}{gray}{0.83}
\definecolor{dark-grey}{gray}{0.3}
\definecolor{text-grey}{gray}{.08}

\DeclareRobustCommand{\ebseries}{\fontseries{eb}\selectfont}
\DeclareTextFontCommand{\texteb}{\ebseries}

% custom underilne
\usepackage{contour}
\usepackage[normalem]{ulem}
\renewcommand{\ULdepth}{1.8pt}
\contourlength{0.8pt}
\newcommand{\myuline}[1]{%
  \uline{\phantom{#1}}%
  \llap{\contour{white}{#1}}%
}


% custom font: helvetica-style
\renewcommand*\familydefault{\sfdefault}


\pagestyle{fancy}
\fancyhf{} % clear all header and footer fields
\fancyfoot{}
\renewcommand{\headrulewidth}{0pt}
\renewcommand{\footrulewidth}{0pt}

% Adjust margins
\addtolength{\oddsidemargin}{-0.5in}
\addtolength{\evensidemargin}{0in}
\addtolength{\textwidth}{1in}
\addtolength{\topmargin}{-.5in}
\addtolength{\textheight}{1.0in}

\urlstyle{same}

\raggedbottom
\raggedright
\setlength{\tabcolsep}{0in}
\setlength{\footskip}{4.1pt}

% Sections formatting - serif
% \titleformat{\section}{
%   \vspace{2pt} \scshape \raggedright\large % header section
% }{}{0em}{}[\color{black} \titlerule \vspace{-5pt}]

% TODO EBSERIES
% sans serif sections
\titleformat {\section}{
    \bfseries \vspace{2pt} \raggedright \large % header section
}{}{0em}{}[\color{light-grey} {\titlerule[2pt]} \vspace{-4pt}]

% only for pdflatex
% Ensure that generate pdf is machine readable/ATS parsable
% \pdfgentounicode=1

%-------------------------
% Custom commands
\newcommand{\resumeItem}[1]{
  \item\small{
    {#1 \vspace{-1pt}}
  }
}

\newcommand{\resumeSubheading}[4]{
  \vspace{-1pt}\item
    \begin{tabular*}{\textwidth}[t]{l@{\extracolsep{\fill}}r}
      \textbf{#1} & {\color{dark-grey}\small #2}\vspace{1pt}\\ % top row of resume entry
      \textit{#3} & {\color{dark-grey} \small #4}\\ % second row of resume entry
    \end{tabular*}\vspace{-4pt}
}

\newcommand{\resumeSubSubheading}[2]{
    \item
    \begin{tabular*}{\textwidth}{l@{\extracolsep{\fill}}r}
      \textit{\small#1} & \textit{\small #2} \\
    \end{tabular*}\vspace{-7pt}
}

\newcommand{\resumeProjectHeading}[2]{
    \item
    \begin{tabular*}{\textwidth}{l@{\extracolsep{\fill}}r}
      #1 & {\color{dark-grey}\small #2} \\
    \end{tabular*}\vspace{-4pt}
}

\newcommand{\resumeSubItem}[1]{\resumeItem{#1}\vspace{-4pt}}

\renewcommand\labelitemii{$\vcenter{\hbox{\tiny$\bullet$}}$}

% CHANGED default leftmargin  0.15 in
\newcommand{\resumeSubHeadingListStart}{\begin{itemize}[leftmargin=0in, label={}]}
\newcommand{\resumeSubHeadingListEnd}{\end{itemize}}
\newcommand{\resumeItemListStart}{\begin{itemize}}
\newcommand{\resumeItemListEnd}{\end{itemize}\vspace{0pt}}

\color{text-grey}

%-------------------------------------------
%%%%%%  RESUME STARTS HERE  %%%%%%%%%%%%%%%%%%%%%%%%%%%%


\begin{document}

%----------HEADING----------
\begin{center}
    \textbf{\Large Очкин Никита Валерьевич} \\ \vspace{5pt}
    \small \texttt{+7 977 194-80-50} \hspace{1pt} $|$
    \hspace{1pt} \href{mailto:nkt.ochkin@gmail.com}{\texttt{nkt.ochkin@gmail.com}} \hspace{1pt} $|$
    \hspace{1pt} \href{https://github.com/namenick91}{\texttt{github.com/namenick91}} \hspace{1pt} $|$
    \hspace{1pt} \texttt{Москва}
    \\ \vspace{-3pt}
\end{center}

%-----------PROFILE-----------
\section{ПРОФИЛЬ}
 \begin{itemize}[leftmargin=0in, label={}]
    \small{\item{
     Студент МГТУ им. Н. Э. Баумана на направлении <<Математика и компьютерные науки>>. Фокус -- прикладная финансовая аналитика: прогнозирование временных рядов, скоринговые модели, клиентская аналитика и обработка данных на Python/SQL.
    }}
 \end{itemize}

%-----------EDUCATION-----------
\section{ОБРАЗОВАНИЕ}
  \resumeSubHeadingListStart
    \resumeSubheading
      {МГТУ им. Н. Э. Баумана}{2022 -- 2026}
      {Бакалавриат, факультет <<Фундаментальные науки>>}{Москва}
        \resumeItemListStart
          \resumeItem{Направление: <<Математика и компьютерные науки>>, профиль: <<Суперкомпьютерное моделирование и искусственный интеллект в инженерных задачах>>}
          \resumeItem{Текущий средний балл: \textbf{4.31}}
        \resumeItemListEnd
  \resumeSubHeadingListEnd


%-----------RESEARCH-----------
\section{НАУЧНАЯ РАБОТА}
  \resumeSubHeadingListStart
    \resumeProjectHeading
      {\textbf{Convformer}: прогнозирование многомерных временных рядов}{2026}
      \resumeItemListStart
        \resumeItem{Препринт по модели прогнозирования временных рядов готовится к публикации в журнале 2--3 квартиля. Препринт: \href{https://github.com/namenick91/Convformer}{github.com/namenick91/Convformer}}
        \resumeItem{Работал с прогнозом на длинных горизонтах для многомерных рядов: модель учитывает локальные изменения, тренд, сезонность и долгосрочные зависимости}
        \resumeItem{На стандартных наборах данных получил более низкую среднеквадратичную ошибку (MSE) на длинных горизонтах без резкого роста вычислительных затрат}
        \resumeItem{Работал с данными которые часто встречаются в финансах: котировки, доходности, макропоказатели и другие ряды, где важны горизонт прогноза и устойчивость ошибки}
      \resumeItemListEnd
  \resumeSubHeadingListEnd


%-----------PROJECTS-----------
\section{ПРОЕКТЫ}
    \resumeSubHeadingListStart
      \resumeProjectHeading
          {\textbf{Customer Churn}: клиентская аналитика и скоринг}{}
          \resumeItemListStart
            \resumeItem{Построил модель оценки вероятности оттока клиента на табличных данных: подготовка данных, построение признаков, кросс-валидация и подбор параметров}
            \resumeItem{Постановка близка к банковским задачам retention и скоринга: вероятностный score, ранжирование клиентов по риску и выбор рабочего порога}
            \resumeItem{Код: \href{https://github.com/namenick91/ML-Workshelf\#customer-churn-prediction}{github.com/namenick91/ML-Workshelf}}
          \resumeItemListEnd
    \resumeSubHeadingListEnd


%-----------PROGRAMMING SKILLS-----------
\section{НАВЫКИ}
 \begin{itemize}[leftmargin=0in, label={}]
    \small{\item{
     \textbf{Финансовая аналитика}{: временные ряды, churn / scoring, антифрод, эконометрика, извлечение данных из документов}\vspace{2pt} \\
     \textbf{Модели и методы}{: градиентный бустинг, регрессия, нейросети, бинарная классификация, оценка качества моделей}\vspace{2pt} \\
     \textbf{Инструменты}{: Python, SQL}
    }}
 \end{itemize}


%-------------------------------------------
\end{document}
